\section{Related work}
\label{sec:related}

\paragraph{Tokenization premiums.} Tokenizers trained mostly on English encode other languages in more tokens, and since commercial models bill per token, those languages pay more for the same content~\citep{petrov2023unfairness,ahia-etal-2023-languages}. Recent work documents such premiums for African, European, Indic and other languages~\citep{lundin-etal-2026-token,ovcharov2026european,somide2026african,tamang2024evaluating}, including a Nepali/English ratio of $4.79\times$ for \texttt{cl100k} on 150 parallel passages~\citep{kadariya2026nepali}. \citet{srivastava2026tokenizertax} and \citet{mandarapu2026ledger} measure large Indic premiums on FLORES-200, and the former reports that \texttt{o200k} cuts the mean Indic premium from $8.0\times$ to $2.1\times$; neither attributes the change to the regex. We report the premium as a sentence-level token ratio on parallel FLORES-200 text~\citep{goyal-etal-2022-flores,nllb2022}, following \citet{petrov2023unfairness}, and add paired bootstrap intervals. Fewer tokens do not by themselves make a better model~\citep{schmidt-etal-2024-tokenization}, which is why we test the repair with continued pretraining.

\paragraph{Pre-tokenization of combining marks.} That GPT-2-style regexes split Indic words at vowel signs was reported publicly in 2024, in a tiktoken issue that also proposed the word class \texttt{[\textbackslash p\{L\}\textbackslash p\{M\}]+}~\citep{tiktoken_issue292}. \citet{minixhofer2024zett} adjusted their GPT-2-based regex so that it does not over-segment languages with combining marks inside words, naming Hindi and Tamil. \citet{velayuthan-sarveswaran-2025-egalitarian} showed that the GPT-2, GPT-4 and Llama-3 pre-tokenizers break Tamil, Sinhala and Hindi text, and found in Tamil training experiments that pre-tokenization largely determines compression. \citet{rana2025mutant} gained 38--40\% in fertility by swapping the GPT-2 regex for Llama-4's. \citet{regmi2026vowel} identified the cause precisely (combining marks fall outside \pL{}), stated the pre-token lower bound, which they call the pre-tokenization ceiling, and measured it across 26 languages. With matched tokenizers and 268M-parameter models trained from scratch, they showed that the wider class lowers Nepali bits per byte by 4.4\%. Among the 1{,}345 Hugging Face text-generation repositories they classified, 63\% use the letters-only class, and these carry 72.5\% of 30-day downloads. They also note that the repair itself is prior art, since \texttt{o200k} ships it. \citet{land2026smallbugs} documents the same bug, lists the families still affected, and reports that Qwen3.5 fixed it. The edit itself is known (tiktoken \#292, ZeTT, \texttt{o200k}, Qwen3.5); what is not studied is repairing a deployed model.

\paragraph{Beyond one character class.} Other work changes more than one class. Grapheme- and akshara-aware pre-tokenizers enforce orthographic syllable boundaries~\citep{kapila2026typedriven,darshana2026wwho}, and grapheme-level tokenizers run subword learning over grapheme clusters~\citep{krishnan2026tamil,velayuthan-sarveswaran-2025-egalitarian}. \citet{shrestha2026papaya} add a morphology-guided pre-tokenizer for Nepali on top of the mark-aware class. MorphTok~\citep{brahma2025morphtok} works at the merge level, constraining BPE so that dependent vowels do not stand alone. SuperBPE and BoundlessBPE~\citep{liu2025superbpe,schmidt2025boundless} lift pre-token boundaries in a second training stage, on the same logic that pre-token boundaries cap compression. All of this work trains new tokenizers. We study the setting where the model already exists.

\paragraph{Adapting the vocabulary of a pretrained model.} Vocabulary extension adds tokens, initialises their embeddings from existing ones~\citep{gee-etal-2022-fast,minixhofer-etal-2022-wechsel,dobler-de-melo-2023-focus,liu-etal-2024-ofa,mundra-etal-2024-empirical} and continues pretraining~\citep{tejaswi-etal-2024-exploring,kim2024eeve,csaki2024sambalingo}. On byte-level models, continued BPE appends new merges to the existing merge list~\citep{purason-etal-2026-teaching,smith2026inplace}. \citet{yamaguchi2026lowres} extend Llama-3 and Qwen models for several abugidas, and their released Telugu tokenizer keeps the Llama-3 regex (\S\ref{sec:ceiling}). These studies cover our model families and several abugidas, but keep the original pre-tokenizer, which caps them as we show in \S\ref{sec:ceiling}. For Nepali, \citet{duwal2024nepalidapt} continue pretraining Llama-3-8B without changing the vocabulary, a setting like our \RO{}. Zero-shot tokenizer transfer~\citep{minixhofer2024zett} applies its mark-aware regex to target tokenizers for SentencePiece models, which have no regex to repair. \citet{dagan2024getting} change the pre-tokenizer of Code Llama during fine-tuning and find that the model recovers its performance after more than 50B tokens of training; their change alters English and code tokenization, which ours does not.

The closest prior system is TituLLM~\citep{nahin-etal-2025-titullms}, which extends Llama-3.2 for Bangla. Its released tokenizer replaces the Llama-3 regex with a \texttt{\textbackslash w}-based one that keeps vowel signs attached, but the paper does not identify the regex as a cause and has no regex-only or vocabulary-only arm. Its regex also changes digit, identifier and whitespace pre-tokens in English. On Nepali its premium ($2.60\times$) is slightly above Llama-3's ($2.58\times$; \S\ref{sec:audit}): a mark-aware regex without merges for the target language does not lower the premium by itself. We give a controlled attribution, and a repair that leaves every input without marks or Devanagari code points unchanged.
